% !TeX root = ../navier-stokes-manuscript.tex
\subsection{What temporal ascent exposes in space}\label{sub:TheViscosityDrivenSpatialAscent}

The recurrence \eqref{eq:BodyTemporalVelocityRecurrence} contains a Laplacian at every temporal level.
That observation says that higher spatial derivatives are present.
It does not yet separate them from transport and pressure.
The critical height supplies a scalar readout of the same response.

For \(s\geq0\), define
\[
  E_s(t):=\frac12\norm{\Lambda^su(t)}_2^2
\]
and the pressure-completed nonlinear transfer
\[
  \mathcal P_s(t)
  :=-
  \operatorname{Re}
  \pair{\Lambda^su(t)}
  {\Lambda^s\bigl((u\cdot\nabla)u+\nabla p\bigr)(t)}_{L^2}.
\]
The gradient part vanishes in this global scalar pairing because \(\Lambda^su\) is divergence-free.
It is kept inside \(\mathcal P_s\) because the vector recurrence remains pressure-complete before and after the contraction.
Pairing the momentum equation at height \(s\) gives
\begin{equation}\label{eq:BodyFullSpatialEnergyBalance}
  E_s'=\mathcal P_s-2\nu E_{s+1}.
\end{equation}
One time derivative of a spatial energy exposes the velocity energy one full Sobolev step higher and retains the signed transfer that produced it.

At the critical height \(H=E_{1/2}\), the first row is
\begin{equation}\label{eq:BodyFirstCriticalRow}
  H'=\mathcal P_{1/2}-2\nu E_{3/2}.
\end{equation}
This identity gives every first-passage interval an exact bill.
Using \(H(t_{k+1})-H(t_k)=L_{k+1}-L_k=L_k\),
\begin{equation}\label{eq:EachZenoPassageCost}
  \int_{t_k}^{t_{k+1}}\mathcal P_{1/2}(t)\,dt
  =L_k+2\nu\int_{t_k}^{t_{k+1}}E_{3/2}(t)\,dt.
\end{equation}
The nonlinear response must supply both the next critical-height gain and the viscous cost accumulated during the same shrinking interval.
This identity does not impose a lower bound on the length of that interval.
The transfer may intensify while the first-passage times compress.

Differentiate \eqref{eq:BodyFirstCriticalRow} once and use \eqref{eq:BodyFullSpatialEnergyBalance} at \(s=3/2\):
\begin{equation}\label{eq:BodySecondCriticalRow}
  H''
  =\partial_t\mathcal P_{1/2}
  -2\nu\mathcal P_{3/2}
  +(2\nu)^2E_{5/2}.
\end{equation}
The next derivative gives
\begin{equation}\label{eq:BodyThirdCriticalRow}
  H'''
  =\partial_t^2\mathcal P_{1/2}
  -2\nu\partial_t\mathcal P_{3/2}
  +(2\nu)^2\mathcal P_{5/2}
  -(2\nu)^3E_{7/2}.
\end{equation}
Every new time derivative reaches one higher spatial energy and preserves every differentiated transfer encountered on the way.
The pattern is triangular.

Induction gives, for \(n\geq1\),
\begin{equation}\label{eq:BodyCriticalHeightSpatialAscent}
  H^{(n)}
  =(-2\nu)^nE_{n+1/2}
  +\sum_{j=0}^{n-1}
  (-2\nu)^j
  \partial_t^{\,n-1-j}\mathcal P_{j+1/2}.
\end{equation}
Indeed, differentiating the row at order \(n\) and replacing
\[
  E_{n+1/2}'
  =\mathcal P_{n+1/2}-2\nu E_{n+3/2}
\]
creates the next transfer term and the next highest spatial energy with the required coefficient.

The temporal velocity row \(V_n\) and the scalar derivative \(H^{(n)}\) have the same index but are different objects.
Their exact relation is
\begin{equation}\label{eq:BodyHeightDerivativeVelocityContraction}
  H^{(n)}
  =\frac12\sum_{a=0}^{n}\binom{n}{a}
  \operatorname{Re}
  \pair{\Lambda^{1/2}V_a}{\Lambda^{1/2}V_{n-a}}_{L^2}.
\end{equation}
The scalar height derivative contracts every temporal velocity row up to level \(n\).
It is neither \(V_n\) nor a norm of \(V_n\).

Keeping the full transfer sum and solving \eqref{eq:BodyCriticalHeightSpatialAscent} for the highest spatial energy gives the completed row
\begin{equation}\label{eq:BodyCompletedCriticalRow}
  \mathscr C_n
  :=
  \frac{
    H^{(n)}
    -\displaystyle\sum_{j=0}^{n-1}
    (-2\nu)^j
    \partial_t^{\,n-1-j}\mathcal P_{j+1/2}
  }{(-2\nu)^n}
  =E_{n+1/2},
  \qquad n\geq1,
\end{equation}
with \(\mathscr C_0:=H=E_{1/2}\).
At index \(n\), the equation has produced three simultaneous readings:
\[
  V_n=\partial_t^nu,
  \qquad
  H^{(n)},
  \qquad
  \mathscr C_n=E_{n+1/2}.
\]
The first is a vector-valued Eulerian response.
The second is a scalar contraction of the temporal rows.
The third is the higher spatial energy isolated after the complete differentiated transfer has been retained.

This is the temporal--spatial reversal.
A proposed singularity would have to lose control at fixed temporal and spatial levels, while differentiating its critical response exposes stronger spatial energies one after another.
The completed-row identities do not bound those energies.
They show what a successful estimate must retain.
The scalar diagonal cannot hold the full spatial column inside each \(V_n\), so temporal depth and spatial depth must now be kept as independent coordinates of the same equation-generated object.
